\documentclass[journal,onecolumn]{IEEEtran}
\usepackage{cite}
\usepackage{amsmath,amssymb,amsfonts}
\usepackage{algorithmic}
\usepackage{graphicx}
\usepackage{textcomp}
\usepackage{xcolor}

\begin{document}

\title{Título del Artículo en Formato IEEE Open Journal}

\author{Nombre del Autor, \IEEEmembership{Member,~IEEE}, y Segundo Autor, \IEEEmembership{Fellow,~IEEE}
\thanks{Este trabajo fue apoyado por...}
\thanks{Los autores pertenecen a la Universidad de... (correo electrónico: autor@ejemplo.com).}}

\maketitle

\begin{abstract}
Este es el resumen del artículo en formato IEEE Open Journal a una sola columna. Debe sintetizar brevemente el objetivo, la metodología, los resultados principales y las conclusiones de la investigación.
\end{abstract}

\begin{IEEEkeywords}
Inteligencia Artificial, IEEE Open Journal, Formato de publicación, Plantilla LaTeX.
\end{IEEEkeywords}

\section{Introducción}
\IEEEPARSTART{E}{ste} documento sirve como plantilla de inicio para la preparación de artículos según las directrices de publicación a una sola columna de la serie de revistas \textit{IEEE Open Journal}.

\section{Metodología}
Describa aquí la metodología utilizada en su trabajo.

\section{Resultados}
Presente los resultados obtenidos con tablas y figuras explicativas.

\section{Conclusiones}
Resuma las principales aportaciones y el trabajo futuro.

\bibliographystyle{IEEEtran}
\bibliography{referencias}

\end{document}
